%%%%%%%%%%%%%%%%%%%%%%%%%%%%%%%%%%%%%%%%%%%%%%%%%%%%%%%%%%%%%%%%%%%%%%%%%%%
% Plantilla genérica para TFG/TFM/Tesis.
%
% Función: configuración del documento y comandos personalizados.
%
% Uso: debes modificar este fichero para adaptarlo a tu trabajo.
% Define la titulación, el título, tus datos, los tutores, el idioma,
% la estructura y las opciones que quieras utilizar.
% Consulta los valores admitidos en los comentarios de cada variable.
% Al final puedes añadir tus propios comandos.
%
% Estos datos se reutilizan en el documento, las páginas institucionales
% y los formularios, para evitar que tengas que introducirlos varias veces.
%%%%%%%%%%%%%%%%%%%%%%%%%%%%%%%%%%%%%%%%%%%%%%%%%%%%%%%%%%%%%%%%%%%%%%%%%%%

%%%%%%%%%%%%%%%%%%%%%%%%%%%%%%%%%%%%%%%%%%%%%%%%%%%%%%%%%%%%%%%%%%%%%%%%%%% 
% INICIO Configuración de las variantes del documento

% Selecciona el tipo de documento (TFM, TFG, tesis, etc.).
% La titulación puede tomar los siguientes valores:
% Los identificadores ajenos a la UAH usan <degree>-<school>-<university>; conserva los guiones.
% Grados de la UAH:
%  En la Escuela Politécnica Superior:
%   GITT   - Grado en Ingeniería en Tecnologías de la Telecomunicación
%   GIEC   - Grado en Ingeniería Electrónica de Comunicaciones
%   GIT    - Grado en Ingeniería Telemática
%   GIST   - Grado en Ingeniería en Sistemas de Telecomunicación
%   GIC    - Grado en Ingeniería de Computadores
%   GII    - Grado en Ingeniería Informática
%   GMC    - Grado en Matemáticas y Computación
%   GISI   - Grado en Ingeniería en Sistemas de Información
%   GIEAI  - Grado en Ingeniería en Electrónica y Automática Industrial
%   GITI   - Grado en Ingeniería en Tecnologías Industriales
%  En la Facultad de Ciencias:
%   GFIE   - Grado en Física e Instrumentación Espacial
%   GBS    - Grado en Biología Sanitaria
%   GB     - Grado en Biología
%   GCA    - Grado en Ciencias Ambientales
%   GQ     - Grado en Química
%   GCCTF  - Grado en Criminalística: Ciencias y Tecnologías Forenses
%
% Másteres de la UAH:
%  En la Escuela Politécnica Superior:
%   MUSEA  - Máster Universitario en Sistemas Electrónicos Avanzados. Sistemas Inteligentes
%   MUIT   - Máster Universitario en Ingeniería de Telecomunicación
%   MUII   - Máster Universitario en Ingeniería Industrial
%   MUIE   - Máster Universitario en Ingeniería Electrónica
%   MUCTE  - Máster Universitario en Ciencia y Tecnología desde el Espacio
%   MUC    - Máster Universitario en Ciberseguridad
%   MUANBD - Máster Universitario en Analítica de Negocio y Big Data
%
% Programas de doctorado:
%  PHDUAH  - Doctorado UAH
%  PHD-ETSIT-UPM - Doctorado de la ETSIT-UPM
%  PHDUAH-COMPENDIUM - Doctorado UAH por compendio de publicaciones
%  PHD-ETSIT-UPM-COMPENDIUM - Doctorado ETSIT-UPM por compendio de publicaciones
%
% Titulaciones de la ESCET-URJC:
%  ITI-ESCET-URJC - Grado en Ingeniería en Tecnologías Industriales (ESCET-URJC)
%
% Titulaciones de la ETSII-URJC (TFG individuales):
%  GCID-ETSII-URJC - Grado en Ciencia e Ingeniería de Datos
%  GDDV-ETSII-URJC - Grado en Diseño y Desarrollo de Videojuegos
%  GII-ETSII-URJC - Grado en Ingeniería Informática
%  GICO-ETSII-URJC - Grado en Ingeniería de Computadores
%  GICI-ETSII-URJC - Grado en Ingeniería de la Ciberseguridad
%  GIS-ETSII-URJC - Grado en Ingeniería del Software
%  GIA-ETSII-URJC - Grado en Inteligencia Artificial
%  GMAT-ETSII-URJC - Grado en Matemáticas
%
% Titulaciones de la EIF-URJC (TFG individuales):
%  GISAM-EIF-URJC - Grado en Ingeniería en Sistemas Audiovisuales y Multimedia
%  GIST-EIF-URJC - Grado en Ingeniería en Sistemas de Telecomunicación
%  GIT-EIF-URJC - Grado en Ingeniería en Telemática
%  GITT-EIF-URJC - Grado en Ingeniería en Tecnologías de Telecomunicación
%  GIAA-EIF-URJC - Grado en Ingeniería Aeroespacial en Aeronavegación
%  GIATA-EIF-URJC - Grado en Ingeniería Aeroespacial en Transporte y Aeropuertos
%  GIAVA-EIF-URJC - Grado en Ingeniería Aeroespacial en Vehículos Aeroespaciales
%  GIA-EIF-URJC - Grado en Ingeniería Aeroespacial
%  GIB-EIF-URJC - Grado en Ingeniería Biomédica
%  GIRS-EIF-URJC - Grado en Ingeniería de Robótica Software
%  GCID-EIF-URJC - Grado en Ciencia e Ingeniería de Datos
%  GP-EIF-URJC - Grado en Paisajismo
%  GFA-EIF-URJC - Grado en Fundamentos de la Arquitectura
%
% Titulaciones de la ETSIT-UPM:
%  GISTS-ETSIT-UPM  - Grado en Ingeniería de Tecnologías y Servicios de Telecomunicación
%  GIB-ETSIT-UPM    - Grado en Ingeniería Biomédica
%  GISD-ETSIT-UPM   - Grado en Ingeniería y Sistemas de Datos
%  MUIT-ETSIT-UPM   - Máster Universitario en Ingeniería de Telecomunicación
%  MUCS-ETSIT-UPM   - Máster Universitario en Ciberseguridad
%  MUESFV-ETSIT-UPM - Máster Universitario en Energía Solar Fotovoltaica
%  MUIB-ETSIT-UPM   - Máster Universitario en Ingeniería Biomédica
%  MUIRST-ETSIT-UPM - Máster Universitario en Ingeniería de Redes y Servicios Telemáticos
%  MUISE-ETSIT-UPM  - Máster Universitario en Ingeniería de Sistemas Electrónicos
%  MUTSC-ETSIT-UPM  - Máster Universitario en Teoría de la Señal y Comunicaciones
%  MUNE-ETSIT-UPM   - Máster Universitario en Neurotecnología (páginas en español o inglés)
%
% Otros documentos:
%  GEINTRARR - Geintra Research Report (soporte preliminar)
%
% Especifica aquí tu titulación, de entre las listadas arriba.
\newcommand{\myDegree}{GIEC}

\newcommand{\mySpecialty}{} % Grados de la UAH y perfil/especialidad en la hoja de control del anteproyecto de TFM EPS-UAH.

% Selecciona el idioma del documento y sus ajustes específicos.
% Los valores admitidos son english o spanish.
\newcommand{\myLanguage}{spanish}

% Estilo visual del documento principal; no afecta al papeleo:
% standard  - aspecto original de la plantilla (opción predeterminada recomendada)
% editorial - diseño con serifas y un separador vertical discreto en los capítulos
% modern    - títulos sin serifas, color para resaltar elementos y números de capítulo grandes
% framed    - títulos de capítulo enmarcados, al estilo Lenny
% shaded    - títulos de capítulo sombreados, al estilo Bjornstrup
% mimosis   - diseño basado en Garamond e inspirado en la plantilla Mimosis
\newcommand{\myTypesettingStyle}{standard}



%%%%%%%%%%%%%%%%%%%%%%%%%%%%%%%%%%%%%%%%%%%%%%%%%%%%%%%%%%%%%%%%%%%%%%%%%%%
% Información general del documento
\newcommand{\myBookTitleSpanish}{Plantilla unificada para la generación de memorias de TFGs, TFMs y tesis doctorales}
\newcommand{\myBookTitleEnglish}{Unified Template for the Generation of TFGs, TFMs and PhD Thesis}
\newcommand{\myThesisKeywords}{Plantillas de trabajos fin de grado/máster y tesis doctorales, \LaTeX, soporte de español e inglés, generación automática} % (máximo de cinco)
\newcommand{\myThesisKeywordsEnglish}{Bsc., Msc. and PhD. Thesis template, \LaTeX, English/Spanish support, automatic generation} % (máximo de cinco)
\newcommand{\myConfidentialContent}{No} % Puede ser Yes o No; se utiliza
                                        % en el papeleo del MUII desde julio de 2021.

%%%%%%%%%%%%%%%%%%%%%%%%%%%%%%%%%%%%%%%%%%%%%%%%%%%%%%%%%%%%%%%%%%%%%%%%%%%
% Datos del autor
\newcommand{\myAuthorName}{Javier}
\newcommand{\myAuthorSurname}{Macías Guarasa}
\newcommand{\myAuthorFullName}{\myAuthorName{} \myAuthorSurname{}}
\newcommand{\myAuthorGender}{male} 
\newcommand{\myAuthorEmail}{javier.maciasguarasa@uah.es}
\newcommand{\myAuthorDNI}{14991977-W}
% Datos personales para la solicitud de anteproyecto; no siempre son necesarios.
\newcommand{\myAuthorStreet}{C/ Calle de la Calle, 22}
\newcommand{\myAuthorCity}{Meco}
\newcommand{\myAuthorPostalCode}{28880}
\newcommand{\myAuthorProvince}{Madrid}
\newcommand{\myAuthorTelephone}{666666666}


%%%%%%%%%%%%%%%%%%%%%%%%%%%%%%%%%%%%%%%%%%%%%%%%%%%%%%%%%%%%%%%%%%%%%%%%%%%
% Datos del tutor
% GISD-ETSIT-UPM: el tutor definido en myAcademicTutorFullName se muestra como
% supervisor técnico, no como supervisor académico.
\newcommand{\myAcademicTutorFullName}{Roberto Barra Chicote} 
\newcommand{\myAcademicTutorGender}{male}
\newcommand{\myAcademicTutorDNI}{20052011-J}
\newcommand{\myAcademicTutorEmail}{roberto.barra@somewhere.com}
% Departamento e institución son campos separados; deja vacío lo que no proceda.
\newcommand{\myAcademicTutorDepartment}{Departamento de Electrónica}
\newcommand{\myAcademicTutorDepartmentEnglish}{Department of Electronics}
\newcommand{\myAcademicTutorInstitution}{Amazon Text-To-Speech Research}
\newcommand{\myAcademicTutorInstitutionEnglish}{Amazon Text-To-Speech Research}
\newcommand{\myAcademicTutorPosition}{Principal Scientist}
\newcommand{\myAcademicTutorPositionEnglish}{Principal Scientist} % Sólo para titulaciones UPM: créditos de tesis en inglés.
\newcommand{\myAcademicTutorForeignUniversity}{Foreign University name}
\newcommand{\myAcademicTutorForeignCity}{Foreign city name}
\newcommand{\myAcademicTutorIsPhD}{true} % true o false, según la normativa de TFM de la EPS-UAH.

% Sólo para titulaciones UPM: Ponente ETSIT-UPM (GISTS/GIB y másteres con su
% mismo modelo): papel distinto del tutor y del cotutor. Déjalo vacío si todavía
% no está asignado; la hoja del tribunal deja el valor vacío.
\newcommand{\myPonenteFullName}{}

% Sólo para titulaciones UPM: Supervisión técnica del MUNE-ETSIT-UPM, distinta de la
% supervisión académica (myAcademicTutorFullName). Puede ser la misma persona;
% indica ambos nombres. Si todavía no está asignada, déjalo vacío para mostrar
% una línea en blanco.
\newcommand{\myUPMTechnicalSupervisorFullName}{}
% GISD-ETSIT-UPM no utiliza myUPMTechnicalSupervisorFullName: usa el tutor y cotutor.

%%%%%%%%%%%%%%%%%%%%%%%%%%%%%%%%%%%%%%%%%%%%%%%%%%%%%%%%%%%%%%%%%%%%%%%%%%%
% Datos del cotutor
% GISD-ETSIT-UPM: el cotutor definido en myCoTutorFullName se muestra como supervisor
% académico; completa su nombre en el campo siguiente.
\newcommand{\myCoTutorFullName}{María Isidra de Guzmán y de la Cerda} % Déjalo vacío si no hay cotutor.
\newcommand{\myCoTutorGender}{female}
\newcommand{\myCoTutorDNI}{17850606-F}
\newcommand{\myCoTutorEmail}{maria.guzman@uah.es}
\newcommand{\myCoTutorDepartment}{Departamento de Filosofía y Letras Humanas}
\newcommand{\myCoTutorDepartmentEnglish}{Department of Philosophy and Humanities}
\newcommand{\myCoTutorInstitution}{Universidad de Alcalá}
\newcommand{\myCoTutorInstitutionEnglish}{University of Alcalá}
\newcommand{\myCoTutorPosition}{Catedrática de Universidad}
\newcommand{\myCoTutorPositionEnglish}{Full Professor} % Sólo para titulaciones UPM: créditos de tesis en inglés.
\newcommand{\myCoTutorIsPhD}{true} % true o false

%%%%%%%%%%%%%%%%%%%%%%%%%%%%%%%%%%%%%%%%%%%%%%%%%%%%%%%%%%%%%%%%%%%%%%%%%%%
% Datos institucionales
% La universidad, su acrónimo y el centro se seleccionan automáticamente en
% Config/degrees.tex y Config/institutions.tex según \myDegree.

% Departamento donde se realiza el trabajo, independiente de las afiliaciones
% del tutor y del cotutor. Deja vacío lo que no proceda: no se infiere ni traduce.
\newcommand{\myDepartment}{Departamento de Electrónica}
\newcommand{\myDepartmentEnglish}{Department of Electronics}

\newcommand{\myPhDProgram}{Programa de Doctorado en Electrónica: Sistemas Electrónicos Avanzados. Sistemas Inteligentes}
\newcommand{\myPhDProgramEnglish}{PhD. Program in Electronics: Advanced Electronic Systems. Intelligent Systems}

% Sólo para titulaciones UPM: Tesis PHD-ETSIT-UPM: utiliza el nombre oficial
% completo del programa en los campos anteriores (Doctorado en ... / Doctoral
% Degree in ...), sin abreviarlo. Titulación previa de acceso al doctorado, no
% el título de doctor que solicitas.
\newcommand{\myUPMdegree}{Ingeniero de Telecomunicación}
\newcommand{\myUPMdegreeEnglish}{Telecommunication Engineer}
% Año de solicitud del depósito; puede ser distinto del año de defensa.
\newcommand{\myThesisDepositYear}{\myThesisDefenseYear}
% Texto completo de financiación competitiva para los créditos PHD-ETSIT-UPM.
% Déjalo vacío si no procede; admite texto LaTeX y no aparece en otras titulaciones.
\newcommand{\myThesisFundingStatement}{}
\newcommand{\myThesisFundingStatementEnglish}{}

\newcommand{\myResearchGroup}{GEINTRA}


%%%%%%%%%%%%%%%%%%%%%%%%%%%%%%%%%%%%%%%%%%%%%%%%%%%%%%%%%%%%%%%%%%%%%%%%%%%
% Miembros del tribunal y personal del departamento
\newcommand{\myTribunalPresident}{Francisca de Nebrija}
\newcommand{\myTribunalPresidentGender}{female}
\newcommand{\myTribunalPresidentDNI}{15141522-R} % Necesario para el compromiso de confidencialidad.
% ETSIT-UPM: el primer vocal (FirstSpokesperson) se muestra como Vocal.
\newcommand{\myTribunalFirstSpokesperson}{Francisco Vallés de Covarrubias}
\newcommand{\myTribunalFirstSpokespersonGender}{male}
\newcommand{\myTribunalFirstSpokespersonDNI}{15241592-K} % Necesario para el compromiso de confidencialidad.
\newcommand{\myTribunalSecondSpokesperson}{Sira Elena Palazuelos Cagigas}
\newcommand{\myTribunalSecondSpokespersonGender}{female}
\newcommand{\myTribunalSecondSpokespersonDNI}{20052007-D} % Necesario para el compromiso de confidencialidad.
\newcommand{\myTribunalAlternateMember}{Elio Antonio de Nebrija}
\newcommand{\myTribunalAlternateMemberGender}{male}
\newcommand{\myTribunalAlternateMemberDNI}{14441522-Y} % Necesario para el compromiso de confidencialidad.
\newcommand{\myTribunalSecretary}{María Castellano Arroyo}
\newcommand{\myTribunalSecretaryGender}{female} % male o female; independiente de la secretaría del departamento
\newcommand{\myDepartmentSecretary}{Departmento de Electrónica} % Para el papeleo de TFG y TFM, incluido MUSEA.
\newcommand{\myDepartmentSecretaryGender}{female}                 % Para el papeleo de TFG y TFM, incluido MUSEA.
\newcommand{\mySchoolDeputyDirectorFullName}{Francisco Jiménez de Cisneros} % Nombre de quien firma por la UAH en los convenios de cooperación educativa de TFG y TFM. Si está vacío, deja la línea del nombre en blanco.
\newcommand{\mySchoolDeputyDirectorGender}{male} % male o female; selecciona Subdirector o Subdirectora de la escuela en esos convenios.
\newcommand{\myTFMComisionPresident}{Felipe Espinosa Zapata} % Nombre utilizado en el papeleo del MUIE.
\newcommand{\myTFMComisionPresidentGender}{male} % male o female; cargo en la hoja de control del anteproyecto de TFM EPS-UAH, independiente del tribunal.

%%%%%%%%%%%%%%%%%%%%%%%%%%%%%%%%%%%%%%%%%%%%%%%%%%%%%%%%%%%%%%%%%%%%%%%%%%%
% Sólo para titulaciones UPM: Tribunal ampliado y revisores externos para tesis
% PHD-ETSIT-UPM. Completa los nombres asignados; deja vacíos los campos que no
% procedan.
\newcommand{\myTribunalThirdSpokesperson}{Francisca de Pedraza}
\newcommand{\myTribunalSecondAlternateMember}{Fernando de Mena}
\newcommand{\myFirstExternalReviewerFullName}{}
\newcommand{\myFirstExternalReviewerInstitution}{Universidad de Alcalá}
\newcommand{\myFirstExternalReviewerInstitutionEnglish}{University of Alcalá}
\newcommand{\mySecondExternalReviewerFullName}{}
\newcommand{\mySecondExternalReviewerInstitution}{Universidad de Alcalá}
\newcommand{\mySecondExternalReviewerInstitutionEnglish}{University of Alcalá}

%%%%%%%%%%%%%%%%%%%%%%%%%%%%%%%%%%%%%%%%%%%%%%%%%%%%%%%%%%%%%%%%%%%%%%%%%%%
% Calificaciones numéricas de 0 a 10, no aportaciones ya ponderadas.
% TFG EPS-UAH: nota final = 0.7 * tutor + 0.3 * tribunal.
% Un suspenso de cualquiera de las partes es vinculante. Si la media diera
% aprobado, la plantilla detiene la compilación para que consultes al departamento.
% TFM: la nota final es la del tribunal; la nota del tutor no interviene.
\newcommand{\myAcademicTutorGrade}{5} % Nota del tutor/cotutor del TFG, sobre 10.
\newcommand{\myTribunalGrade}{10} % Nota del tribunal del TFG o nota final del TFM, sobre 10.
\newcommand{\myTribunalMHProposal}{NO} % Los valores admitidos son YES o NO.


%%%%%%%%%%%%%%%%%%%%%%%%%%%%%%%%%%%%%%%%%%%%%%%%%%%%%%%%%%%%%%%%%%%%%%%%%%%
% Fechas
\newcommand{\myThesisProposalDate}{16 de enero de 2027} % Fecha del anteproyecto.

\newcommand{\myThesisDepositDate}{15 de junio de 2027}
\newcommand{\myThesisDepositDateEnglish}{June 15\textsuperscript{th}, 2027}

% En los RR, myThesisDefenseDate es la fecha que aparece en la portada.
\newcommand{\myThesisAcademicYear}{2026/2027}
\newcommand{\myThesisDefenseYear}{2027}
\newcommand{\myThesisDefenseDate}{24 de junio de \myThesisDefenseYear{}}
\newcommand{\myThesisdefenseDateEnglish}{June 24\textsuperscript{th}, \myThesisDefenseYear{}}
% Si prefieres el formato británico para la fecha en inglés, usa este ejemplo:
% \newcommand{\myThesisdefenseDateEnglish}{6\textsuperscript{th} of January, 2018}

\newcommand{\myPaperworkDate}{\myThesisDepositDate}
\newcommand{\myPaperworkDateEnglish}{\myThesisDepositDateEnglish}

%%%%%%%%%%%%%%%%%%%%%%%%%%%%%%%%%%%%%%%%%%%%%%%%%%%%%%%%%%%%%%%%%%%%%%%%%%%
% Datos para la publicación en abierto
\newcommand{\myAuthorizationOpenPublishing}{Yes}
\newcommand{\myAuthorizationOpenPublishingEmbargoMonths}{0} % Puede ser
                                                            %  0,  6, 12,
                                                            % 18, 24

\newcommand{\myResearchVicerrector}{Excmo. Sr. Francisco J. de la Mata de la Mata}
\newcommand{\myResearchVicerrectorGender}{male}

% Derechos de autor
\newcommand{\myCopyrightStatement}{\myAuthorFullName{}. Some rights reserved. This document is under terms of Creative Commons license Attribution - Non Commercial - Non Derivatives.}
\newcommand{\myLicenseURL}{http://creativecommons.org/licenses/by-nc-nd/3.0/es/}

\newcommand{\myResearchReportID}{RR-2027-01}

%%%%%%%%%%%%%%%%%%%%%%%%%%%%%%%%%%%%%%%%%%%%%%%%%%%%%%%%%%%%%%%%%%%%%%%%%%%
% Elementos opcionales del documento. Cada valor debe ser exactamente true o false.
% Desactivar una lista no desactiva la función LaTeX correspondiente.
\newcommand{\myIncludeSampleLetter}{false}
\newcommand{\myIncludeDedication}{true}
\newcommand{\myIncludeAcknowledgements}{true}
\newcommand{\myIncludeListOfFigures}{true}
\newcommand{\myIncludeListOfTables}{true}
\newcommand{\myIncludeListOfAcronyms}{true}
\newcommand{\myIncludeListOfSymbols}{true}
\newcommand{\myIncludeListOfCode}{true}
\newcommand{\myIncludeListOfAlgorithms}{true}
\newcommand{\myIncludeListOfVideos}{true}

% Bibliografía del libro: puedes incluirla por capítulo, general o ambas.
% Los dos primeros valores deben ser true o false.
\newcommand{\myIncludeChapterBibliographies}{false}
\newcommand{\myIncludeFinalBibliography}{true}
% global: conserva el número de cada referencia en todo el documento.
% chapter: empieza en 1 en cada capítulo. La bibliografía general, si se
% incluye, tiene su propia numeración independiente.
\newcommand{\myBibliographyNumbering}{global}

% En las páginas que admiten ambas formas, true separa tutor y cotutor.
% false muestra una denominación conjunta, adaptada al tipo de trabajo.
% No altera las páginas de formato institucional fijo (por ejemplo, EIF).
% Consulta el alcance de esta opción en el manual.
\newcommand{\myFlagSplittedAdvisors}{true}

%%%%%%%%%%%%%%%%%%%%%%%%%%%%%%%%%%%%%%%%%%%%%%%%%%%%%%%%%%%%%%%%%%%%%%%%%%%
% Colores de los enlaces
% Enlaces de los índices general, de tablas y de figuras.
%\newcommand{\mytoclinkcolor}{blue}
\newcommand{\mytoclinkcolor}{black}
%\newcommand{\myloflinkcolor}{red}
\newcommand{\myloflinkcolor}{black}
%\newcommand{\mylotlinkcolor}{green}
\newcommand{\mylotlinkcolor}{black}

% Enlaces de las listas de código, algoritmos y vídeos.
%\newcommand{\myothertoclinkcolor}{magenta}
\newcommand{\myothertoclinkcolor}{black}

% Otros enlaces internos del documento.
\newcommand{\mylinkcolor}{blue}
%\newcommand{\mylinkcolor}{black}

% Enlaces a URL y citas bibliográficas.
\newcommand{\myurlcolor}{blue}
%\newcommand{\myurlcolor}{black}
%\newcommand{\mycitecolor}{green}
\newcommand{\mycitecolor}{blue}
%\newcommand{\mycitecolor}{black}

% FIN Configuración de las variantes del documento
%%%%%%%%%%%%%%%%%%%%%%%%%%%%%%%%%%%%%%%%%%%%%%%%%%%%%%%%%%%%%%%%%%%%%%%%%%% 

%%%%%%%%%%%%%%%%%%%%%%%%%%%%%%%%%%%%%%%%%%%%%%%%%%%%%%%%%%%%%%%%%%%%%%%%%%% 
% INICIO Comandos personalizados
% Define aquí tus propios comandos.

% Este comando define el formato del texto en inglés dentro de un documento
% en español.
\DeclareRobustCommand{\texten}[1]{\textit{#1}}

\def\ci{\perp\!\!\!\perp}

% Ejemplos de comandos de uso frecuente.
\newcommand{\circulo}{\large $\circ$}
\newcommand{\asterisco}{$\ast$}
\newcommand{\cuadrado}{\tiny $\square$}
\newcommand{\triangulo}{\scriptsize $\vartriangle$}
\newcommand{\triangv}{\scriptsize $\triangledown$}
\newcommand{\diamante}{\large $\diamond$}

\newcommand{\new}[1]{\textcolor{magenta}{#1 }}
\newcommand{\argmax}[1]{\underset{#1}{\operatorname{argmax}}}

% Estos comandos se utilizan en los capítulos de ejemplo.
\newcommand{\verticalSpacingSRPMaps}{-0.3cm}

\newcommand{\CAVDB}{\texttt{CAV3D}}
\newcommand{\AVThreeTVLoc}{\texttt{AV3T-VLoc}}
\newcommand{\LandPoseVLoc}{\texttt{LandPose-VLoc}}
\newcommand{\DeltaPoseAVT}{$\Delta_{\AVThreeTVLoc}^{\LandPoseVLoc}$}
\newcommand{\mae}{\epsilon{}}
\newcommand{\maeThreeD}{\mae_{3D}}
\newcommand{\maeThreeDfine}{\maeThreeD'}
\newcommand{\maeTwoD}{\mae_{2D}}
\newcommand{\maeTwoDfine}{\maeTwoD'}

\newcommand{\NumDataFrames}{N_F}


% Anotaciones de revisión
\newcommand{\todojmg}{\todo[author=JMG, inline, color=green!20, caption={}]}
\newcommand{\todorbc}{\todo[author=RBC, inline, color=red!20,   caption={}]}
\newcommand{\todomig}{\todo[author=MIG, inline, color=blue!20,  caption={}]}
% \newcommand{\todospc}{\todo[author=SPC, inline, color=green!60,  caption={}]}
% \newcommand{\todolbm}{\todo[author=LBM, inline, color=green!40, caption={}]}
% \newcommand{\todojmb}{\todo[author=JMB, inline, color=red!40,   caption={}]}
% \newcommand{\todoapt}{\todo[author=APT, inline, color=blue!40,  caption={}]}

% FIN Comandos personalizados
%%%%%%%%%%%%%%%%%%%%%%%%%%%%%%%%%%%%%%%%%%%%%%%%%%%%%%%%%%%%%%%%%%%%%%%%%%% 


%%%%%%%%%%%%%%%%%%%%%%%%%%%%%%%%%%%%%%%%%%%%%%%%%%%%%%%%%%%%%%%%%%%%%%%%%%%
% INICIO Definiciones obsoletas o de menor utilidad; consérvalas por compatibilidad.
\newcommand{\myFirstAdvisorFullName}{\myAcademicTutorFullName} % Obsoleta: toma el nombre del tutor.
\newcommand{\mySecondAdvisorFullName}{\myCoTutorFullName} % Obsoleta: toma el nombre del cotutor.
\newcommand{\myFirstAdvisorDNI}{\myAcademicTutorDNI} % Obsoleta: toma el DNI del tutor.
\newcommand{\mySecondAdvisorDNI}{\myCoTutorDNI} % Obsoleta: toma el DNI del cotutor.
\newcommand{\mybookFigure}{alumno} % Obsoleta: indicaba en los TFG
                                   % la condición del autor que firmaba
                                   % el convenio (normalmente "alumno").

% FIN Definiciones obsoletas o de menor utilidad
%%%%%%%%%%%%%%%%%%%%%%%%%%%%%%%%%%%%%%%%%%%%%%%%%%%%%%%%%%%%%%%%%%%%%%%%%%%

%%% Local Variables:
%%% TeX-master: "../book"
%%% End:
